\section{Formal Analysis}

\label{sec:analysis}

\subsection{Self-Improvement as Optimization}

We formalize the self-improvement objective in terms of tool success rates.

\begin{definition}[Tool Success Rate]
\label{def:success_rate}
The success rate $\mu_t^{(s)}$ of tool $t$ at session $s$ is:
\begin{equation}
\mu_t^{(s)} = \E\bigl[\mathbf{1}[\tau.\text{outcome}=\textsc{success}]
               \mid \tau.\text{tool}=t,\; \Sigma^{(s)}\bigr],
\end{equation}
where the expectation is over the task distribution conditional on the agent
state at session $s$.
\end{definition}

\begin{definition}[Self-Improvement Objective]
\label{def:objective}
The self-improvement objective is:
\begin{equation}
\max_{\{\mathcal{K}^{(s)}, \mathcal{F}^{(s)}\}_{s \geq 1}} \lim_{S \to \infty}
\frac{1}{S} \sum_{s=1}^{S} \E\biggl[\frac{1}{|\mathcal{V}_T|} \sum_{t \in \mathcal{V}_T} \mu_t^{(s)}\biggr],
\end{equation}
i.e., maximize the long-run average tool success rate over all tools.
\end{definition}

\subsection{Convergence of the 4-Loop System}

\begin{theorem}[4-Loop Convergence]
\label{thm:convergence}
Under the following conditions:
\begin{enumerate}[leftmargin=1.5em, nosep]
    \item[\textbf{(A1)}] The task distribution $\mathcal{D}$ is stationary.
    \item[\textbf{(A2)}] The BIN build pipeline succeeds with probability $p_b > 0$
          on any triggered $(t,f)$ pair.
    \item[\textbf{(A3)}] The space of failure modes $\mathcal{V}_F$ is finite.
    \item[\textbf{(A4)}] The agent does not regress: adding a new tool does not
          decrease the success rate of existing tools.
\end{enumerate}
the 4-Loop System achieves:
\begin{equation}
\lim_{S \to \infty} \frac{1}{|\mathcal{V}_T|} \sum_{t \in \mathcal{V}_T} \mu_t^{(S)} = 1
\quad \text{almost surely.}
\end{equation}
\end{theorem}

\begin{proof}
We proceed by showing that every $(t,f)$ pair with positive occurrence
probability is eventually eliminated.

\textit{Step 1.} Fix any $(t, f)$ with positive probability $p_{tf} > 0$.
Under stationarity (A1), the expected number of sessions until $\theta = 3$
failures of $(t, f)$ occur in a single session is finite:
$\E[\tau_{tf}] \leq 1 / p_{tf}^\theta < \infty$.

\textit{Step 2.} When the BIN trigger fires for $(t,f)$, the build pipeline
succeeds with probability $p_b > 0$ (A2). If it fails, the same session
continues accumulating failures, so the trigger fires again in the next session
by Step 1. By the second Borel-Cantelli lemma, the build pipeline succeeds
infinitely often, so with probability 1, there exists a session $s^*_{tf}$
after which the $(t,f)$ pair is eliminated from the active failure modes.

\textit{Step 3.} By finiteness of $\mathcal{V}_F$ (A3), the number of distinct
$(t,f)$ pairs is $|\mathcal{V}_T| \times |\mathcal{V}_F| < \infty$. Therefore
$s^* = \max_{t,f} s^*_{tf} < \infty$ with probability 1.

\textit{Step 4.} By the no-regression assumption (A4), the tool library $\mathcal{K}^{(s)}$
is monotonically improving. After session $s^*$, all failure modes have been
addressed and $\mu_t^{(s)} \to 1$ for all $t$ as $s \to \infty$.
\end{proof}

\begin{remark}
Assumption (A4) is the strongest assumption and may fail in practice if new
tools have bugs. The probationary status mechanism in the BIN pipeline
(Section~\ref{sec:loop1}) is designed to mitigate this by quarantining
unproven tools until they accumulate positive evidence.
\end{remark}

\subsection{Information-Theoretic Analysis}
\label{sec:info_theory}

We now show formally that append-only telemetry logs preserve more
decision-relevant information than narrative summaries.

\begin{definition}[Decision-Relevant Information]
\label{def:relevant_info}
Given telemetry log $\mathcal{T}$, the decision-relevant information for the
BIN trigger is:
\begin{equation}
I_{\text{BIN}}(\mathcal{T}) = H\bigl(\{(t,f) : \#_{tf}(\mathcal{T}) \geq \theta\}\bigr),
\end{equation}
where $H$ denotes Shannon entropy and $\#_{tf}(\mathcal{T})$ is the count of
$(t,f)$ failures in $\mathcal{T}$.
\end{definition}

\begin{theorem}[Telemetry Entropy Bound]
\label{thm:entropy}
Let $\mathcal{T}_n$ be a telemetry log of $n$ records and let $\hat{\mathcal{T}}_k$ be
any narrative summary of $\mathcal{T}_n$ using $k$ tokens ($k \ll n$). Then:
\begin{equation}
I_{\text{BIN}}(\mathcal{T}_n) \geq I_{\text{BIN}}(\hat{\mathcal{T}}_k) + \Omega\!\left(\log n\right),
\end{equation}
with high probability over the distribution of narrative summaries.
\end{theorem}

\begin{proof}[Proof sketch]
The BIN trigger requires exact counts $\#_{tf}(\mathcal{T}) \geq \theta$ for
specific $(t,f)$ pairs. Reconstructing exact counts from a $k$-token summary
requires at least $\log_2(\max_{t,f} \#_{tf}) \approx \log_2(n)$ bits per
$(t,f)$ pair that is near the threshold. A $k$-token summary has at most
$k \cdot \log_2 |\mathcal{V}|$ bits of total information. For $k \ll n$, this
cannot represent all pairs at threshold-resolution, so information about
near-threshold pairs is discarded. The direct telemetry log represents every
record exactly, preserving $O(\log n)$ more bits of threshold-resolution
information per near-threshold pair.
\end{proof}

\subsection{Deterministic vs.\ LLM-Based Triggers}
\label{sec:triggers}

We analyze the precision and recall of two trigger mechanisms for the BIN loop:
deterministic counting (our approach) vs.\ LLM-based self-judgment.

\begin{definition}[Trigger Precision and Recall]
\label{def:precision_recall}
For a trigger mechanism $M$ and ground truth $G$ (set of $(t,f)$ pairs that
genuinely benefit from a new tool):
\begin{align}
P(M) &= \frac{|M \cap G|}{|M|}, \quad R(M) = \frac{|M \cap G|}{|G|}.
\end{align}
\end{definition}

\begin{proposition}[Deterministic Trigger Dominance]
\label{prop:trigger_dominance}
For $\theta = 3$ and a stationary failure mode distribution, the deterministic
counting trigger $M_D$ satisfies:
\begin{equation}
P(M_D) \geq P(M_{\text{LLM}}) \quad \text{and} \quad R(M_D) \geq R(M_{\text{LLM}}),
\end{equation}
with probability $\geq 1 - \delta$ for any $\delta > 0$ when $n \geq n_0(\delta)$.
\end{proposition}

The empirical support for this proposition comes from our evaluation
(Section~\ref{sec:evaluation}), where we measure $P(M_D) = 0.94$, $R(M_D) = 0.89$
against $P(M_{\text{LLM}}) = 0.61$, $R(M_{\text{LLM}}) = 0.72$. The LLM
trigger has higher recall in low-count regimes (detecting problems with fewer
than $\theta$ occurrences) but much lower precision (many false positives from
LLM hallucination about problems that did not occur frequently).

\begin{table}[H]
\centering
\caption{Threshold sensitivity analysis for BIN trigger at $\theta \in \{2,3,4,5\}$.}
\label{tab:threshold_sensitivity}
\begin{tabular}{ccccc}
\toprule
$\theta$ & Precision & Recall & F1 & Builds/Session \\
\midrule
2 & 0.81 & 0.95 & 0.87 & 2.3 \\
3 & \textbf{0.94} & \textbf{0.89} & \textbf{0.91} & 1.1 \\
4 & 0.97 & 0.72 & 0.83 & 0.6 \\
5 & 0.99 & 0.58 & 0.73 & 0.3 \\
\midrule
LLM & 0.61 & 0.72 & 0.66 & 1.8 \\
\bottomrule
\end{tabular}
\end{table}

Table~\ref{tab:threshold_sensitivity} confirms that $\theta = 3$ is Pareto-optimal:
it achieves the highest F1 score while keeping the builds-per-session rate
below 1.5 (a reasonable upper bound for the 5-minute latency budget per session).

% ============================================================================
% 10. ANTI-PATTERNS AND DESIGN DECISIONS
% ============================================================================
